% LSPV-1 REPORT::HIERFIRE1::2026-10-12T18:11:39Z
\documentclass[11pt]{article}
\usepackage{amsmath,amssymb,amsthm,booktabs,enumitem,geometry,hyperref,longtable}
\geometry{margin=1in}
\hypersetup{colorlinks=true,linkcolor=blue,urlcolor=blue}
\newtheorem{theorem}{Theorem}
\newtheorem{proposition}{Proposition}
\newcommand{\OPEN}{\textsc{open}}
\newcommand{\EXACT}{\textsc{proved/exact}}
\newcommand{\KNOWN}{\textsc{known theorem}}
\newcommand{\RET}{\textsc{falsified/retired}}
\newcommand{\COMP}{\textsc{computational}}
\title{LSPV-1: An Encoding-Scale Firewall for Generic Time-Hierarchy Transfers}
\author{UNIVERSAL\_CORE\_LOWER\_BOUND}
\date{2026-10-13 03:11 KST}

\begin{document}
\maketitle

\paragraph{Executive status.}
\textbf{P versus NP remains open.}  This run proves an accounting theorem, not a
SAT lower bound.  A standard computation-to-SAT-to-clean-core pipeline does not
turn the mere assumption ``the clean core is in polynomial time'' into a
deterministic time-hierarchy contradiction: the known simulation has essentially
linear (indeed $T\log T$) length in a $T$-step computation, while the assumed
solver may have an arbitrary fixed exponent.  The missing arrow remains a direct,
model-matched lower bound for a fixed NP-complete language, or a genuine
sublinear-in-$T$ external contraction that works for the \emph{same realization}
and survives all construction and access costs.

\section{Predecessor recovery and corpus reviewed}
\paragraph{Predecessor write-recovery preflight.}
The newest canonical LSPV-5 artifact was read back as
\texttt{research/main/2026-10-13\_0148\_MAIN\_TRACECONJ1.tex}; its required MAIN
header and final sentinel are present and its body is substantive.  Hence no
recovery request or duplicate file was warranted.

\paragraph{CORPUS REVIEWED.}
The default-branch tree and current root \texttt{main.tex} were surveyed before
substantive work.  I read \texttt{AGENTS.md}, \texttt{CODEX.md}, the complete
\texttt{.codex/skills/pnp-research/SKILL.md} protocol and mode files, all current
task specifications under \texttt{tasks/}, the newest verified MAIN, and the
newest verified LSPV-1 through LSPV-4 reports.  Material historical items included
LSPV-1 \texttt{CUTGAP1}, \texttt{INFONOGO1}, and \texttt{LB-0018}; LSPV-2
\texttt{INFOFLOWCHAIN1}; LSPV-3 \texttt{REVUNCOMP1}; LSPV-4
\texttt{READCEIL1}; and the cost-certificate spine.  I also inspected the actual
\texttt{main4.tex} and \texttt{main5.tex} entry points, governance/write notes,
and the HEC, VSC-HEC, G-PPC-HEC, CRS-PPC, and CCSPC exploration branch.  Repository
search covered ``time hierarchy'', ``Cook--Levin'', ``encoded length'',
``decision tree'', ``evasive'', ``padding'', ``succinct'', ``wrapper'', and
``same realization''.

\paragraph{Primary-literature preflight.}
The two blocking obstacles are (i) the arbitrary polynomial exponent of a
hypothetical SAT decider and (ii) the realized size/construction cost of encoding
a long computation.  Primary sources checked were Cook's original NP-completeness
paper, Hartmanis--Stearns on deterministic time hierarchy, Robson's
$O(T\log T)$ RAM-computation-to-SAT construction, and Fortnow--Lipton--van
Melkebeek--Viglas on time--space lower bounds for SAT.  The decision-tree survey
of Buhrman--de Wolf was also checked to locate, but not re-claim, the older linear
query ceiling.  URLs are recorded in Section~\ref{sec:sources}.

\section{Exact candidate and accounting theorem}
Fix a binary encoding of a language $C\subseteq\{0,1\}^*$ that is NP-complete
under a specified uniform many-one reduction.  Quantifiers below are over ordinary
deterministic uniform random-access or multitape machines with all input,
construction, reduction, and output costs charged.

\begin{theorem}[HIERARCHY-ENCODING-SCALE-FIREWALL1]\label{thm:firewall}
Let $M$ decide a language $L$ on length-$n$ inputs within $T(n)$ steps.  Suppose a
fully materialized reduction pipeline constructs, for every input $x$, a clean-core
instance $R(M,x)$ of encoded length
\[
       N(n)=T(n)^{\lambda+o(1)},\qquad \lambda>0,
\]
in time $T(n)^{\rho+o(1)}$.  Suppose further that a deterministic uniform decider
$A$ for $C$ has worst-case running time $N^{K_A+o(1)}$, where without loss of
generality $K_A\ge 1$.  Then the generic algorithm ``construct $R(M,x)$ and run
$A$'' is certified only by the bound
\[
  T(n)^{\max\{\rho,\lambda K_A\}+o(1)}.
\]
Consequently this bookkeeping yields a strict exponent speedup over the original
$T(n)$ computation only if
\[
           \rho<1 \quad\text{and}\quad \lambda K_A<1.
\]
In particular, a standard same-scale encoding with $\lambda\ge1$ cannot, from the
polynomial-time hypothesis alone, supply such a hierarchy speedup for any
$K_A\ge1$.
\end{theorem}

\begin{proof}
The constructed instance has length $T^{\lambda+o(1)}$, so substituting it into
the stated worst-case bound for $A$ gives
$(T^{\lambda+o(1)})^{K_A+o(1)}=T^{\lambda K_A+o(1)}$.
Sequentially composing construction and solution costs is bounded by the larger
exponent.  A deterministic time-hierarchy contradiction requires an independently
established algorithm whose exponent is strictly below the chosen hierarchy
threshold, giving the two displayed strict inequalities.  If $\lambda\ge1$ and
$K_A\ge1$, the second inequality fails.  This is an upper-bound accounting fact;
it does not assert that every encoding has $\lambda\ge1$, nor that $A$ cannot be
faster on the reduction image.
\end{proof}

Robson's standard RAM simulation produces a Boolean expression of size
$O(T\log T)$, hence $\lambda=1$ at exponent scale.  A subsequent fixed polynomial
reduction of degree $d\ge1$ to a chosen clean core changes this to
$\lambda=d$ (up to subpolynomial factors), worsening rather than repairing the
inequality.  Cook's tableau simulation has the same qualitative firewall.

\paragraph{Exact scope of the negative result.}
The route
\[
  C\in\mathrm P
  \Longrightarrow \text{standard materialized computation encoding}
  \Longrightarrow \text{time-hierarchy contradiction}
\]
is \RET{} as a \emph{generic} argument when its realized encoding scale is
$\lambda\ge1$.  This does not retire direct lower bounds for $C$, reductions with
a proved $\lambda<1/K_A$, or algorithms proved specially fast on the reduction
image.  Those would be new theorems, not consequences of ``$C\in\mathrm P$'' alone.

\section{SAME\_REALIZATION, model matching, and hostile controls}
The simulated $T$-step computation, the materialized formula, and the run of $A$
are three different realizations.  The theorem charges their composition; it does
not transfer a structural formula lower bound to an arbitrary program.  A succinct
circuit describing a tableau is a different input representation and therefore a
different language unless one proves a model-preserving compilation with all
costs.  An oracle call likewise suppresses the very input-access cost that must be
charged in the target uniform machine.

\begin{longtable}{p{0.24\linewidth}p{0.18\linewidth}p{0.48\linewidth}}
\toprule
Control & Status & Effect on this candidate\\
\midrule
2SAT, Horn, affine & exact control & These promise families may be easy; a generic
NP-complete reduction image need not lie in them.  If it does, that is a new
image-specific theorem.\\
Bounded treewidth & exact control & A bounded-width image would permit dynamic
programming, but bounded width must be proved for the same reduction image.\\
Independent products & failed shortcut & Product structure can increase semantic
counts without changing the exponent identity $\lambda K_A$.\\
Padding & \RET{} as repair & Padding replaces the source size by a larger one and
does not make a materialized $T$-step encoding sublinear in $T$.\\
Redundant clauses & exact control & They only enlarge $N$; semantic equivalence is
not an algorithmic contraction.\\
Alternate encodings & \OPEN{} & Relevant only with a proved realized exponent
$\lambda<1/K_A$ and construction/access exponent below one.\\
Wrappers & exact control & Equivalent wrappers can hide internal states but cannot
alter the charged end-to-end exponent of this explicit pipeline.\\
One-bit output & exact control & Decision requires only one bit, so no output-size
argument forces tableau reconstruction.\\
Succinct tableau & failed transfer & This changes the representation/problem unless
a fully charged compilation theorem is supplied.\\
\bottomrule
\end{longtable}

\paragraph{Relativization/algebrization/Natural-Proofs audit.}
The accounting theorem is syntactic and relativizes; therefore it cannot itself
overcome Baker--Gill--Solovay.  It neither uses arithmetization nor evades the
Aaronson--Wigderson algebrization barrier.  It asserts no large constructive circuit
property, so Razborov--Rudich is not directly triggered.  These observations are
limitations, not evidence for a separation.

\paragraph{Structural versus arbitrary computation.}
Tableau size, proof-system length, branching-program width, Schur-square dimension,
and quadratic-relation rank are representation/model facts.  None is a lower bound
for every uniform deterministic decider absent a transfer theorem to the same
charged realization.  Current unconditional SAT time--space lower bounds apply to
restricted resource regimes and do not supply the required all-polynomial-time
contradiction.

\section{Finite falsification audit}
The bounded program
\texttt{experiments/review/hierarchy\_encoding\_scale\_firewall\_audit.py}
checked 4,672 exact arithmetic/control identities with zero failures.  It enumerated
$K=1,\ldots,24$, multiple rational $\lambda$, padding powers $1,\ldots,12$, a
concrete $N(T)=T\lceil\log_2(T+2)\rceil$ scale for $T=2,\ldots,2^{20}$, and ten
hostile controls.  The result is
\texttt{results/raw/review/2026-10-13\_hierarchy\_encoding\_scale\_firewall.json}.
Script SHA-256:
\texttt{cfb2f596432edb722e2c09a4f55863eb9fad6163a2c171adf452317ae02c4010};
result SHA-256:
\texttt{b1d7196c9052ae003376cc8164dac01f271789b88a769e2c903e5254adc05e8b}.
Status: \COMP{}.  This finite audit tests bookkeeping only and is not asymptotic
evidence for P versus NP.

\section{Post-quantum and cross-domain transfer ledger}
\begin{itemize}[leftmargin=*]
\item \textbf{LWE worst-case/average-case architecture: ANALOGY/CANDIDATE needing a
new bridge.}  External black-box extraction is preferable to forcing a solver to
materialize a tableau, but no unconditional NP worst-case/average-case reduction or
sublinear charged contraction is known here.
\item \textbf{Schur squares, quadratic-relation ranks, Hilbert functions, syzygies:
FAILED TRANSFER for this candidate.}  They may distinguish a chosen code or
representation, but supply no same-realization runtime lower bound for arbitrary
deciders.  Any hidden algebraic structure may instead be algorithmically exploitable.
\item \textbf{Lossiness/collision entropy: ANALOGY/CANDIDATE.}  The exact inequality
$|\operatorname{im}L|\ge |A|^2/\operatorname{Coll}(L)$ is valid, but an image-size
bound is not a time bound without an external extraction/charging theorem.
\item \textbf{VLSI area--time and communication cuts: useful structural analogy
requiring a new transfer theorem.}  Their locality/routing assumptions do not hold
for arbitrary Turing/RAM programs.
\item \textbf{Time hierarchy: directly transferable theorem only after a strict
speedup is proved in the same machine model.}  The present firewall identifies why
standard SAT encodings do not provide that speedup.
\end{itemize}

\section{Duplication, contradictions, and reconciliation}
\paragraph{DUPLICATION / CONTRADICTION AUDIT.}
The tempting theorem ``a Boolean function on $N$ bits has decision-tree complexity
at most $N$'' is already present in LSPV-1 \texttt{LB-0018} and sharpened by
LSPV-4 \texttt{READCEIL1}; it was therefore not reissued.  Polynomial query closure
and exogenous-challenge flattening were also already proved.  The cost-certificate
spine's quasipolynomial-window theorem is consistent with the present result but has
a different antecedent.  No prior report located by repository search stated the
explicit $\lambda K_A<1$ firewall for a generic materialized hierarchy reduction.

\paragraph{WHAT CHANGED RELATIVE TO ROOT MAIN.}
The LSPV-1 specialty now records a precise negative theorem for one solver-independent
route: ordinary Cook/Robson-scale encodings do not convert an arbitrary polynomial
SAT exponent into a hierarchy speedup.  This narrows the live direct target to either
(a) a direct lower bound for a fixed NP-complete clean core, or (b) a fully charged
external contraction with realized exponent below $1/K_A$ for the same solver.

\paragraph{ADOPTED / REJECTED / RETIRED PRIOR COMPONENTS.}
Adopted: \texttt{INFONOGO1}'s distinction between information volume and time,
\texttt{READCEIL1}'s universal read-certificate ceiling, and the root requirement
that fixed-core NP-completeness suffices.  Rejected: semantic-state count as runtime,
internal-state exposure, erasure, fixed checkpoint schedules, and proof-size transfer
without acquisition.  Retired here: the generic same-scale computation-to-SAT
hierarchy shortcut.

\paragraph{STALE OR OVERSTATED CLAIMS FOUND.}
Any wording that a linear read/certificate bound is a candidate superpolynomial
obstruction is stale after \texttt{READCEIL1}.  Any wording that time hierarchy plus
Cook--Levin alone contradicts $\mathrm{SAT}\in\mathrm P$ omits the encoded-length and
arbitrary-exponent ledger.  Actual \texttt{main4.tex} and \texttt{main5.tex} do not
prove an arbitrary-algorithm acquisition/coverage theorem; their relevant arrows
remain \OPEN{}.

\section{Status and next target}
\begin{center}
\begin{tabular}{p{0.55\linewidth}p{0.27\linewidth}}
\toprule
Claim & Status\\
\midrule
HIERARCHY-ENCODING-SCALE-FIREWALL1 & \EXACT{} (pending independent review)\\
Robson $O(T\log T)$ RAM-to-SAT encoding & \KNOWN{}\\
Finite exponent/control audit & \COMP{}\\
Generic same-scale Cook/Robson hierarchy shortcut & \RET{}\\
All-encoding lower bound $N\ge T^{1-o(1)}$ & \OPEN{} and not claimed\\
Image-specific fast behavior of an arbitrary core solver & \OPEN{}\\
Direct superpolynomial lower bound for a fixed NP-complete core & \OPEN{}\\
$\mathrm P\ne\mathrm{NP}$ & \OPEN{}\\
\bottomrule
\end{tabular}
\end{center}

The highest-value next target is an \emph{external} black-box contraction theorem
with exact quantifiers:
for every uniform decider $A$ of one fixed NP-complete clean core, construct from
$A$ and a $T$-step computation an equivalent charged instance/interaction of
realized scale $T^{\lambda_A+o(1)}$ with $\lambda_A K_A<1$, including extraction,
materialization, access, and reconstruction.  Without this arrow, time hierarchy
does not close the proof.

\section{Primary sources}\label{sec:sources}
\begin{itemize}[leftmargin=*]
\item S. A. Cook, ``The Complexity of Theorem-Proving Procedures'' (1971),
\url{https://doi.org/10.1145/800157.805047}.
\item J. Hartmanis and R. E. Stearns, ``On the Computational Complexity of
Algorithms'' (1965),
\url{https://www3.risc.jku.at/people/schreine/courses/compcomp/HartmanisStearns1965.pdf}.
\item J. M. Robson, ``An $O(T\log T)$ reduction from RAM computations to
satisfiability'' (1991), \url{https://doi.org/10.1016/0304-3975(91)90177-4}.
\item L. Fortnow, R. Lipton, D. van Melkebeek, and A. Viglas,
``Time-Space Lower Bounds for Satisfiability'',
\url{https://lance.fortnow.com/papers/files/tst.pdf}.
\item H. Buhrman and R. de Wolf, ``Complexity Measures and Decision Tree
Complexity: A Survey'', \url{https://arxiv.org/abs/quant-ph/0206169}.
\end{itemize}

% END_OF_VERIFIED_REPORT::LSPV-1::HIERFIRE1
\end{document}
